\documentclass[12pt,a4paper]{article}
\usepackage[utf8]{inputenc}
\usepackage[T1]{fontenc}
\usepackage[brazil]{babel}
\usepackage{amsmath,amssymb}
\usepackage{geometry}
\geometry{margin=2.5cm}
\title{Material de Estudo: An\'alise Dimensional}
\author{Princ\'ipio de Modelagem Matem\'atica}
\date{Unidade 03}
\begin{document}
\maketitle

\section*{Objetivo}
Usar a an\'alise dimensional para verificar equa\c{c}\~oes, deduzir rela\c{c}\~oes entre vari\'aveis e simplificar modelos sem resolver as equa\c{c}\~oes completas.

\section*{Conceitos Fundamentais}

\subsection*{Dimens\~oes e Unidades}
Dimens\~oes s\~ao categorias f\'isicas fundamentais: comprimento $[L]$, massa $[M]$, tempo $[T]$, temperatura $[\Theta]$, corrente $[I]$.
Unidades s\~ao escolhas de escala (metro, quilograma, segundo, \ldots).

\textbf{Regra de ouro:} equa\c{c}\~oes f\'isicas v\'alidas t\^em \emph{dimens\~oes homog\^eneas} --- s\'o se somam termos de mesma dimens\~ao.

\subsection*{Teorema $\Pi$ de Buckingham}
Se um fen\^omeno depende de $n$ vari\'aveis com $k$ dimens\~oes independentes, ent\~ao pode ser descrito por $n - k$ grupos adimensionais $\Pi_i$.

\textbf{Procedimento:}
\begin{enumerate}
  \item Liste as $n$ vari\'aveis e suas dimens\~oes.
  \item Identifique $k$ vari\'aveis \emph{repetidas} (que cobrem todas as dimens\~oes).
  \item Forme os grupos $\Pi_i$ combinando as repetidas com cada vari\'avel restante.
  \item Determine os expoentes pela condi\c{c}\~ao de adimensionalidade.
\end{enumerate}

\subsection*{N\'umeros Adimensionais Importantes}
\begin{itemize}
  \item \textbf{Reynolds:} $Re = \rho v L / \mu$ \quad (iner\'cia/viscosidade --- regime laminar vs.\ turbulento).
  \item \textbf{Froude:} $Fr = v/\sqrt{gL}$ \quad (iner\'cia/gravidade --- escoamentos com superf\'icie livre).
  \item \textbf{Mach:} $Ma = v/c$ \quad (velocidade/velocidade do som).
  \item \textbf{Prandtl:} $Pr = \mu c_p / \kappa$ \quad (difusividade moment\^anea/t\'ermica).
\end{itemize}

\section*{Exemplos Resolvidos}

\subsection*{Exemplo 1 --- Verifica\c{c}\~ao de equa\c{c}\~ao}
Verifique se $F = mv^2/r$ \'e dimensionalmente correta (for\c{c}a centr\'ipeta).
\[
[F] = \frac{[M][L/T]^2}{[L]} = \frac{[M][L^2/T^2]}{[L]} = [ML/T^2] = [F]\quad\checkmark
\]

\subsection*{Exemplo 2 --- Per\'iodo do p\^endulo simples (Buckingham)}
Vari\'aveis relevantes: per\'iodo $T$ $[T]$, comprimento $\ell$ $[L]$, gravidade $g$ $[L/T^2]$, massa $m$ $[M]$.
\begin{itemize}
  \item $n=4$ vari\'aveis, $k=3$ dimens\~oes ($M, L, T$).
  \item $n-k = 1$ grupo adimensional.
\end{itemize}
Formando $\Pi = T^a\,\ell^b\,g^c\,m^d$:
\[
[T]^a[L]^b[LT^{-2}]^c[M]^d = T^0L^0M^0
\]
\[
M:\; d=0;\quad L:\; b+c=0;\quad T:\; a-2c=0.
\]
Escolhendo $c=-1 \Rightarrow b=1,\; a=2$. Logo $\Pi = T^2 g / \ell$, ou seja:
\[
T = C\sqrt{\ell/g}, \quad C = 2\pi \text{ (determinado experimentalmente)}.
\]
\textbf{An\'alise dimensional deu a forma correta sem resolver nenhuma EDO!}

\subsection*{Exemplo 3 --- C\'alculo de Reynolds}
\'Agua ($\rho=1000\,\text{kg/m}^3$, $\mu=10^{-3}\,\text{Pa}\cdot\text{s}$) em cano de di\^ametro $D=2\,\text{cm}$ com velocidade $v=1\,\text{m/s}$:
\[
Re = \frac{\rho v D}{\mu} = \frac{1000 \times 1 \times 0{,}02}{10^{-3}} = 20\,000.
\]
$Re > 4000 \Rightarrow$ escoamento turbulento.

\section*{Exerc\'icios}

\begin{enumerate}
  \item \textbf{(Verifica\c{c}\~ao)} Verifique a homogeneidade dimensional de cada equa\c{c}\~ao:
    \begin{enumerate}
      \item[$a.$] $E = \frac{1}{2}mv^2$ (energia cin\'etica).
      \item[$b.$] $p = \rho g h$ (press\~ao hidrost\'atica).
      \item[$c.$] $\lambda = h/(mv)$ (comprimento de onda de de Broglie, $[h]=[ML^2/T]$).
    \end{enumerate}
  
  \item \textbf{(Buckingham)} A velocidade de uma onda na superf\'icie de um l\'iquido depende de $\rho$ $[M/L^3]$, $\sigma$ (tens\~ao superficial) $[M/T^2]$ e $\lambda$ (comprimento de onda) $[L]$. Use o teorema $\Pi$ para deduzir a forma da velocidade $v$.
  
  \item \textbf{(Reynolds)} Para qual velocidade m\'inima o escoamento de \'oleo ($\rho=900\,\text{kg/m}^3$, $\mu=0{,}1\,\text{Pa}\cdot\text{s}$) em cano de $D=5\,\text{cm}$ se torna turbulento ($Re > 4000$)?
  
  \item \textbf{(Queda livre)} A velocidade terminal de uma esfera caindo num fluido depende de $g$ $[L/T^2]$, $r$ (raio) $[L]$, $\rho_s$ (densidade da esfera), $\rho_f$ (densidade do fluido) e $\mu$ (viscosidade). Quantos grupos adimensionais existem?
  
  \item \textbf{(Energia nuclear)} A energia liberada por uma explos\~ao nuclear depende de $E$ $[ML^2/T^2]$, $\rho_0$ (densidade do ar) $[M/L^3]$ e $t$ (tempo) $[T]$. G.I.\ Taylor (1950) usou fotos secretas da bomba at\^omica para estimar $E$ a partir do raio $R$ da bola de fogo. Mostre que $R \propto (E/\rho_0)^{1/5}\,t^{2/5}$.
  
  \item \textbf{(Froude)} Um navio de comprimento $L=100\,\text{m}$ viaja a $v=10\,\text{m/s}$. Calcule $Fr$ e interprete o resultado.
  
  \item \textbf{(Interpreta\c{c}\~ao)} Explique por que o n\'umero de Reynolds \'e o mesmo para \'agua num r\'io e ar num t\'unel de vento se as velocidades e dimens\~oes forem ajustadas proporcionalmente. Qual \'e a utilidade pr\'atica disso?
\end{enumerate}

\section*{Resumo R\'apido}
\begin{itemize}
  \item Toda equa\c{c}\~ao f\'isica v\'alida \'e dimensionalmente homog\^enea.
  \item Buckingham: $n$ vari\'aveis, $k$ dimens\~oes $\Rightarrow$ $n-k$ grupos adimensionais.
  \item N\'umeros adimensionais ($Re$, $Fr$, $Ma$) caracterizam regimes f\'isicos.
\end{itemize}

\section*{Refer\^encias}
\begin{itemize}
  \item Barenblatt, G.\ I.\ \textit{Scaling, Self-similarity, and Intermediate Asymptotics}. Cambridge, 1996.
  \item White, F.\ M.\ \textit{Fluid Mechanics}. McGraw-Hill, 2011.
  \item Notas de aula da disciplina.
\end{itemize}

\end{document}
